We study service-menu optimization for many-to-one demand-responsive transit (DRT), where each passenger-facing offer is a bundle of a meeting point and a pickup time window. The display decision is formulated as a Multinomial Logit (MNL) assortment problem with operationally endogenous bundle costs: predicted routing cost, capacity risk, pickup timing, and price jointly determine menu value. The solver uses exact subset enumeration for small candidate sets and greedy forward selection for larger sets, with exact-vs-greedy diagnostics quantifying approximation quality. The empirical program is tiered: the RC (Clustered-Random) benchmark supports exact-gap and robust-filtering mechanism diagnostics, low/medium/high uptake regimes identify when menu design is behaviorally active, and Austin/Seattle provide external directional checks. Filtering based on estimated time of arrival (ETA) can reshape the feasible meeting-point assortment before optimization acts; exact-small/greedy-large menu construction yields a transparent operational approximation whose gap is empirically auditable; and pricing/menu effects are meaningful only in regimes where acceptance and non-home uptake are observable rather than nearly degenerate.
